
% ==============================================================
% Section 7: Conclusion
% Author: Quân (RW)
% File:   paper/sections/07_conclusion.tex
% ==============================================================

\section{Conclusion}\label{sec:conclusion}

This paper evaluated GPT-4o mini as an automated quality gate for
GitHub bug reports, measuring its agreement with an annotator
consensus across the Observed Behavior (OB), Expected Behavior (EB),
and Steps to Reproduce (S2R) dimensions using Cohen's Kappa.
Agreement analysis was conducted on 30 annotated issues from three
OSS repositories (pandas, scikit-learn, VS~Code).

The model showed strong alignment with annotators on OB (100\% raw
agreement; Kappa undefined due to no label variation) and substantial
agreement on EB ($\kappa = 0.695$), both at or above the pre-registered
threshold of $\kappa \geq 0.60$.  Agreement on S2R was substantially
lower ($\kappa = 0.417$), falling below the threshold.  The overall
LLM-vs-annotator Kappa ($\kappa = 0.618$) passed the pilot gate but
did not reach the proposal threshold of $\kappa \geq 0.70$.
Separately, a scale-up run on 250 reports confirmed operational
feasibility: 0 of 250 outputs were invalid and the total cost was
approximately US\$0.086.

These results indicate that GPT-4o mini is suitable as a first-pass
filter for OB and EB quality, but should not be used as a standalone
assessor for S2R without human review.  A hybrid pipeline---where the
model screens OB and EB while annotators handle S2R on borderline
cases---better reflects the gate's demonstrated capability profile.

Future work should evaluate prompt \texttt{v3.0-candidate-s2r-relaxed}
against human labels to determine whether the S2R gap can be closed,
and extend the annotated dataset beyond the current 30-issue subset to
support more robust agreement estimates.
